% Magnon Observables
% Converted on 2026-10-08 from docs/lswt/02-observables/magnon-observables.md (status: draft, source-section: "Physical Quantities in Linear Spin Wave Theory (source p. 12, Table I); Number and Spin moment from Correlation Matrix (source pp. 15–16)").
% Review status: draft; awaiting the user's physics and mathematics review.

\hypertarget{sec-lswt-magnon-observables}{%
\section{Magnon Observables}\label{sec-lswt-magnon-observables}}

After the paraunitary diagonalization of \lswtnoteref{LSWT paraunitary diagonalization}{Paraunitary Diagonalization}, every LSWT observable is built from two ingredients at each momentum of the magnetic Brillouin zone (MBZ): the magnon energies \(\varepsilon_{n\mathbf k}\) and the transformation \(\mathsf T_{\mathbf k}\). This document collects the quantities that follow directly from these ingredients, namely the magnon spectrum, the equal-time correlation matrix of the Holstein--Primakoff bosons, the boson number, and the reduced ordered moment. It also serves as an index to the documents that own the remaining observables.

\subsection{Magnon Spectrum and Band Structure}

The magnon energies \(\varepsilon_{n\mathbf k}\), \(n=1,\ldots,N_{\mathrm{sub}}\), are the positive eigenvalues of \(\Sigma_3\mathsf H_{\mathbf k}\) and are periodic over the MBZ. Each band is a branch \(\varepsilon_{n\mathbf k}\) that continues across momentum; at crossings the band label is a matter of convention, since the energies are ordered at each momentum. The negative eigenvalues \(-\varepsilon_{n,-\mathbf k}\) are the particle--hole partners and carry no additional information.

The magnetic unit cell of an ordered state may contain several crystallographic unit cells. A band structure plotted along a path of the crystallographic Brillouin zone (CBZ) then shows all \(N_{\mathrm{sub}}\) bands of the magnetic cell at every momentum, including branches folded back by magnetic reciprocal-lattice vectors. Which folded branch carries weight at a given external momentum is determined by the one-magnon weights of \lswtnoteref{LSWT spin correlations}{Structure Factor and Spectral Function}, not by the energies alone. At a momentum where \(\mathsf H_{\mathbf k}\) is only positive semidefinite, some \(\varepsilon_{n\mathbf k}\) vanishes; the energies remain well defined as limits, although \(\mathsf T_{\mathbf k}\) may not exist there.

\subsection{Ground-State Energy}

The LSWT ground-state energy \(E_{\mathrm{GS}}=E_{\mathrm{cl}}+\Delta E_{\mathrm{zp}}\) and the zero-point correction \(\Delta E_{\mathrm{zp}}=\frac12\sum_{\mathbf k}(\sum_n\varepsilon_{n\mathbf k}-\operatorname{Tr}\mathsf A_{\mathbf k})\) are derived with the diagonal magnon Hamiltonian in \lswtnoteref{LSWT paraunitary diagonalization}{Paraunitary Diagonalization}. The energy per magnetic site is \(\mathcal E_{\mathrm{GS}}=E_{\mathrm{GS}}/N_{\mathrm{site}}\). The correction depends on the energies alone, not on \(\mathsf T_{\mathbf k}\), and it inherits the provisional trace convention of \lswtnoteref{LSWT boson Hamiltonian}{Momentum-Space BdG Hamiltonian}.

\subsection{Equal-Time Correlation Matrix}

The equal-time averages of all products of two Holstein--Primakoff bosons at momentum \(\mathbf k\) form the \(2N_{\mathrm{sub}}\times2N_{\mathrm{sub}}\) correlation matrix

\begin{equation}\label{eq-lswt-boson-correlation-matrix}
\Big[\big\langle\hat\Psi_{\mathbf k}\hat\Psi_{\mathbf k}^\dagger\big\rangle\Big]
=\begin{pmatrix}
\langle\hat a_{\mathbf k\mu}\hat a_{\mathbf k\nu}^\dagger\rangle & \langle\hat a_{\mathbf k\mu}\hat a_{-\mathbf k\nu}\rangle\\
\langle\hat a_{-\mathbf k\mu}^\dagger\hat a_{\mathbf k\nu}^\dagger\rangle & \langle\hat a_{-\mathbf k\mu}^\dagger\hat a_{-\mathbf k\nu}\rangle
\end{pmatrix}_{\mu\nu}
=\mathsf T_{\mathbf k}\,\mathsf N_{\mathbf k}(0)\,\mathsf T_{\mathbf k}^\dagger,
\end{equation}

where each block is an \(N_{\mathrm{sub}}\times N_{\mathrm{sub}}\) matrix with the sublattice indices \(\mu,\nu\). The second equality substitutes \(\hat\Psi_{\mathbf k}=\mathsf T_{\mathbf k}\hat\Phi_{\mathbf k}\). In the Gibbs state the magnon correlation matrix \(\mathsf N_{\mathbf k}(0)=\operatorname{diag}(1+n_{1\mathbf k},\ldots,n_{1,-\mathbf k},\ldots)\) of \lswtnoteref{LSWT spin correlations}{Spin Correlations} is diagonal, with the occupations \(n_{n\mathbf k}=n_{\mathrm B}(\varepsilon_{n\mathbf k})\) of Thermodynamics (Section~\ref{sec-lswt-thermodynamics}). In the ground state, the magnon vacuum, the occupations vanish and \(\mathsf N_{\mathbf k}(0)=\operatorname{diag}(\mathsf I_{N_{\mathrm{sub}}},0)\). The anomalous averages \(\langle\hat a\hat a\rangle\) are nonzero even at zero temperature when \(\mathsf H_{\mathbf k}\) contains the anomalous block \(\mathsf B_{\mathbf k}\).

\subsection{Boson Number and Reduced Moment}

The average number of Holstein--Primakoff bosons on sublattice \(\mu\) is \(\langle\hat n_\mu\rangle=N_{\mathrm{uc}}^{-1}\sum_i\langle\hat a_{i\mu}^\dagger\hat a_{i\mu}\rangle=N_{\mathrm{uc}}^{-1}\sum_{\mathbf k}\langle\hat a_{\mathbf k\mu}^\dagger\hat a_{\mathbf k\mu}\rangle\). The lower-right block of \eqref{eq-lswt-boson-correlation-matrix} contains these averages at \(-\mathbf k\), and the momentum sum covers all momenta, so

\begin{equation}\label{eq-lswt-boson-number}
\langle\hat n_\mu\rangle
=\frac{1}{N_{\mathrm{uc}}}\sum_{\mathbf k\in\mathrm{MBZ}}\Big[\mathsf T_{\mathbf k}\mathsf N_{\mathbf k}(0)\mathsf T_{\mathbf k}^\dagger\Big]_{N_{\mathrm{sub}}+\mu,\,N_{\mathrm{sub}}+\mu}
=\frac{1}{N_{\mathrm{uc}}}\sum_{\mathbf k\in\mathrm{MBZ}}\sum_{n=1}^{N_{\mathrm{sub}}}
\Big[\big|(\mathsf P_{\mathbf k})_{\mu n}\big|^2n_{n\mathbf k}+\big|(\mathsf Q_{\mathbf k})_{\mu n}\big|^2\big(1+n_{n\mathbf k}\big)\Big].
\end{equation}

The second form uses the blocks of the Bogoliubov transformation \(\hat a_{\mathbf k\mu}=\sum_n[(\mathsf P_{\mathbf k})_{\mu n}\hat b_{\mathbf kn}+(\mathsf Q_{-\mathbf k})_{\mu n}\hat b_{-\mathbf kn}^\dagger]\) of \lswtnoteref{LSWT paraunitary diagonalization}{Paraunitary Diagonalization} and relabels \(-\mathbf k\to\mathbf k\) in the hole-column terms. The \(\mathsf Q\) term survives at zero temperature: the quantum reduction \(N_{\mathrm{uc}}^{-1}\sum_{\mathbf k,n}|(\mathsf Q_{\mathbf k})_{\mu n}|^2\) arises only from the mixing of creation and annihilation operators and vanishes when \(\mathsf B_{\mathbf k}=0\), as for a collinear ferromagnet whose Hamiltonian conserves the total spin component along the ordered moment. Exchange anisotropy transverse to the moment, such as \(J^{xx}\ne J^{yy}\) in the local frame or Kitaev and off-diagonal couplings, produces \(\mathsf B_{\mathbf k}\ne0\) and a quantum reduction even in a collinear ferromagnet. The \(\mathsf P\) and \(\mathsf Q\) terms proportional to \(n_{n\mathbf k}\) give the thermal reduction.

The boson number reduces the ordered moment. With \(\hat{\widetilde S}_I^0=S_I-\hat n_I\) and the local longitudinal direction \(\mathbf n_\mu\) of \lswtnoteref{LSWT classical order}{Classical Order and Local Frame}, the average spin of sublattice \(\mu\) is

\begin{equation}\label{eq-lswt-reduced-moment}
\mathbf m_\mu=\big(S_\mu-\langle\hat n_\mu\rangle\big)\,\mathbf n_\mu,
\end{equation}

since the transverse components are linear in the bosons and average to zero in the Gibbs state of the quadratic Hamiltonian. Equation~\eqref{eq-lswt-reduced-moment} gives the length of the moment to order \(S^0\). It gives the full moment to that order only when symmetry fixes the classical spin directions, for example a collinear state along an axis that the Hamiltonian preserves. For a canted or otherwise noncollinear reference configuration whose angles \(\boldsymbol\theta\) are not fixed by symmetry, the zero-point energy \(\Delta E_{\mathrm{zp}}(\boldsymbol\theta)\) of order \(S\) shifts the angles that minimize \(E_{\mathrm{cl}}+\Delta E_{\mathrm{zp}}\) by
\[
\delta\boldsymbol\theta=-\big(\partial_{\boldsymbol\theta}^2E_{\mathrm{cl}}\big)^{-1}\partial_{\boldsymbol\theta}\Delta E_{\mathrm{zp}},
\]
which is of order \(S^{-1}\) because \(E_{\mathrm{cl}}\) is of order \(S^2\). The resulting rotation of the moment, \(S_\mu\,(\partial\mathbf n_\mu/\partial\boldsymbol\theta)\cdot\delta\boldsymbol\theta\), is of order \(S^0\), the same order as \(\langle\hat n_\mu\rangle\). In the boson language it is the transverse expectation value of order \(S^0\) generated by the cubic terms of the Holstein--Primakoff expansion, which the quadratic Hamiltonian omits. To order \(S^0\) the moment of such a state is therefore
\[
\mathbf m_\mu=\big(S_\mu-\langle\hat n_\mu\rangle\big)\,\mathbf n_\mu+S_\mu\,\frac{\partial\mathbf n_\mu}{\partial\boldsymbol\theta}\cdot\delta\boldsymbol\theta,
\]
and the magnetization obtained from \(-\mathrm d(E_{\mathrm{cl}}+\Delta E_{\mathrm{zp}})/\mathrm dh\), with \(h\) the strength of a uniform Zeeman field, contains both terms. The magnetization per site is \(N_{\mathrm{sub}}^{-1}\sum_\mu\mathbf m_\mu\). These moments enter the elastic Bragg scattering of \lswtnoteref{LSWT spin correlations}{Spin Correlations}. The reduction \(\langle\hat n_\mu\rangle\) is of order \(S^0\), compared with \(S_\mu\), so the expansion is controlled only while \(\langle\hat n_\mu\rangle\ll S_\mu\). A value \(\langle\hat n_\mu\rangle\ge S_\mu\) would remove or reverse the moment and signals that LSWT about the chosen reference configuration no longer applies.

\subsubsection{Zero Modes}

At a zero mode the summand of \eqref{eq-lswt-boson-number} can diverge. Near a Goldstone mode of an antiferromagnet, \(\varepsilon_{n\mathbf k}\propto|\mathbf k-\mathbf k_0|\) and \(|\mathsf P_{\mathbf k}|^2,|\mathsf Q_{\mathbf k}|^2\propto1/\varepsilon_{n\mathbf k}\), as for the single mode of \lswtnoteref{LSWT paraunitary diagonalization}{Paraunitary Diagonalization}. At zero temperature the summand then grows as \(1/|\mathbf k-\mathbf k_0|\), which is integrable in two dimensions, and the quantum reduction is finite. At nonzero temperature \(n_{\mathrm B}(\varepsilon)\simeq k_{\mathrm B}T/\varepsilon\) adds a factor \(1/\varepsilon\), and the summand grows as \(|\mathbf k-\mathbf k_0|^{-2}\); the same power follows for a ferromagnetic Goldstone mode with \(\varepsilon\propto|\mathbf k-\mathbf k_0|^2\) and bounded \(\mathsf P_{\mathbf k}\). In two dimensions this gives a logarithmically divergent thermal boson number. The divergence is the LSWT form of the absence of spontaneous continuous-symmetry breaking at nonzero temperature in two dimensions, discussed in \lswtnoteref{LSWT spin Hamiltonian}{LSWT Overview}. The thermodynamic potentials of Thermodynamics (Section~\ref{sec-lswt-thermodynamics}) remain finite in the same situation.

\subsection{Index of Observables}

\begin{longtable}[]{@{}
  >{\raggedright\arraybackslash}p{(\columnwidth - 4\tabcolsep) * \real{0.3333}}
  >{\raggedright\arraybackslash}p{(\columnwidth - 4\tabcolsep) * \real{0.3333}}
  >{\raggedright\arraybackslash}p{(\columnwidth - 4\tabcolsep) * \real{0.3333}}@{}}
\toprule\noalign{}
\begin{minipage}[b]{\linewidth}\raggedright
Observable
\end{minipage} & \begin{minipage}[b]{\linewidth}\raggedright
Ingredients
\end{minipage} & \begin{minipage}[b]{\linewidth}\raggedright
Owner document
\end{minipage} \\
\midrule\noalign{}
\endhead
\bottomrule\noalign{}
\endlastfoot
Magnon energies and bands & \(\varepsilon_{n\mathbf k}\) & This document; \lswtnoteref{LSWT paraunitary diagonalization}{Paraunitary Diagonalization} \\
Ground-state energy, zero-point correction & \(\varepsilon_{n\mathbf k}\), \(\operatorname{Tr}\mathsf A_{\mathbf k}\) & \lswtnoteref{LSWT paraunitary diagonalization}{Paraunitary Diagonalization} \\
Boson number, reduced moment, magnetization & \(\mathsf T_{\mathbf k}\), \(n_{n\mathbf k}\) & This document \\
Partition function, \(F\), \(U\), \(\mathcal S\), \(C\) & \(\varepsilon_{n\mathbf k}\) & Thermodynamics (Section~\ref{sec-lswt-thermodynamics}) \\
Spin correlation function, elastic correlation & \(\mathsf T_{\mathbf q}\), \(\mathsf N_{\mathbf q}(t)\), \(\mathbf m_\mu\) & \lswtnoteref{LSWT spin correlations}{Spin Correlations} \\
Static and dynamic structure factors, spectral function & One-magnon weights \(W_n^{\alpha\beta}\) & \lswtnoteref{LSWT spin correlations}{Structure Factor and Spectral Function} \\
Berry curvature, Chern number, thermal Hall conductivity & \(\mathsf T_{\mathbf k}\), \(\partial_{\mathbf k}\mathsf H_{\mathbf k}\) & \lswtnoteref{LSWT magnon topology}{Topological Magnon Quantities} \\
Classical lower bound and candidate orders & Classical energy only & \lswtnoteref{LSWT classical order}{Appendix: Luttinger–Tisza Method} \\
\end{longtable}

\subsection{References}

\subsubsection{Internal Documents}

\begin{itemize}
\tightlist
\item
  \lswtnoteref{LSWT paraunitary diagonalization}{Paraunitary Diagonalization}: supplies \(\varepsilon_{n\mathbf k}\), \(\mathsf T_{\mathbf k}\) with its blocks \(\mathsf P_{\mathbf k}\) and \(\mathsf Q_{\mathbf k}\), and the ground-state energy.
\item
  \lswtnoteref{LSWT boson Hamiltonian}{Momentum-Space BdG Hamiltonian}: defines the Nambu spinor and the BdG blocks.
\item
  \lswtnoteref{LSWT classical order}{Classical Order and Local Frame}: defines the ordered directions \(\mathbf n_\mu\).
\item
  Thermodynamics (Section~\ref{sec-lswt-thermodynamics}): defines the occupations and the thermodynamic potentials.
\item
  \lswtnoteref{LSWT spin correlations}{Spin Correlations}: defines \(\mathsf N_{\mathbf q}(t)\) and uses the reduced moments.
\item
  \lswtnoteref{LSWT spin correlations}{Structure Factor and Spectral Function}: determines the spectral weight of folded bands.
\item
  \lswtnoteref{LSWT magnon topology}{Topological Magnon Quantities}: builds topological responses from the bands and eigenvectors.
\item
  \lswtnoteref{LSWT classical order}{Appendix: Luttinger–Tisza Method}: bounds the classical energy of candidate reference configurations.
\item
  \lswtnoteref{LSWT spin Hamiltonian}{LSWT Overview}: states the finite-temperature qualification in two dimensions.
\end{itemize}

\subsubsection{External Sources}

\begin{itemize}
\tightlist
\item
  None.
\end{itemize}
